\begin{definition}[Lewis's planar polynomial rigidity, as a hypothesis on $p$]
Let $p\in\mathbb{R}$.  We say that \emph{$p$ satisfies Lewis rigidity} when the following
holds: for every integer $m\ge2$, every continuous $m$-linear form
$Q\colon\mathrm{Plane}^m\to\mathbb{R}$ and every $G\colon\mathrm{Plane}\to\mathrm{Plane}$, if
the degree-$m$ homogeneous polynomial
\[
  P(x):=Q[x,\ldots,x]\qquad(x\in\mathrm{Plane}),
\]
obtained by evaluating $Q$ on $m$ copies of $x$, is $p$-harmonic on all of $\mathrm{Plane}$
with weak gradient $G$, then $P\equiv0$, i.e.\ $Q[x,\ldots,x]=0$ for every
$x\in\mathrm{Plane}$.

This is the planar polynomial rigidity theorem of Lewis (\emph{Smoothness of certain
degenerate elliptic equations}, Proc. Amer. Math. Soc. \textbf{80} (1980), 259--265; see also
the introduction of Lewis--Vogel, \emph{On $p$ Laplace polynomial solutions}, J. Anal.
\textbf{24} (2016), 143--166), restricted to homogeneous polynomials, which is the form in
which the paper invokes it.  In the paper's proof of the vanishing of higher derivatives at a
critical point, the polynomial in question is
$P_m(x)=\tfrac1{m!}D^mu(x_0)[x^m]$, which is exactly of the above shape with
$Q=\tfrac1{m!}D^mu(x_0)$ the $m$-th derivative tensor of $u$ at $x_0$, a continuous $m$-linear
form on $\mathrm{Plane}$.

The hypothesis is stated as a property of $p$ rather than proved: it is one of the results
the paper cites from the literature.
\end{definition}
